% DOCX paragraph 003.
\begin{abstract}
Complex geoscientific tasks often require multiple tools to work together and exchange intermediate results across computational steps. Existing methods commonly plan around specific tools, tying workflows to particular implementations, whereas reusing entire strategies makes it difficult to combine local steps from different strategies according to task requirements. We present AgentGS, a method for capability composition and tool orchestration over skill hypergraphs. AgentGS uses a three-layer Skill--Capability--Tool representation to decouple task strategies, reusable operations, and concrete tools, allowing it to organize an operational workflow before selecting its implementations. The skill hypergraph connects strategies through shared capabilities and provides a structured basis for composition across strategies. Given selected strategies, the system expands their required operations and prerequisite steps. When the current operation chain lacks necessary intermediate steps, it searches for feasible connecting branches along dependencies declared by other strategies. By jointly checking step dependencies, data compatibility, and tool availability, the system incorporates branches that satisfy the constraints without ambiguity, thereby reusing local steps from other strategies. It then selects concrete tools, propagates task data and intermediate results, and locally recovers failed steps. AgentGS can also organize newly composed capability workflows into candidate Skills for subsequent registration and reuse. On GeoToolEval, which contains 603 evaluation records across ten geoscientific domains, AgentGS achieves a mean Final Answer Score (FAS) of 0.8731. Ablation experiments show that hierarchical planning and missing-capability completion both improve task performance, with prerequisite expansion, capability completion, and execution recovery being particularly helpful for multistep tasks.
\end{abstract}
